% Glosario de terminos del trabajo. Se imprime en orden alfabetico por el campo sort.
% Para agregar uno:
% \newglossaryentry{clave}{name={Termino}, sort={termino}, description={Definicion.}}
% El campo sort va sin tildes para que el orden sea correcto.

\newglossaryentry{ajustefino}{name={Ajuste fino (\textit{fine-tuning})}, sort={ajuste fino},
  description={Reentrenamiento parcial de las capas finales de una red preentrenada para adaptarla a una tarea concreta.}}
\newglossaryentry{aleteo}{name={Aleteo de manos}, sort={aleteo de manos},
  description={Movimiento repetitivo y r\'itmico de las manos o los brazos, similar al batir de alas. Es una de las tres conductas que identifica el sistema.}}
\newglossaryentry{aprendizajeautomatico}{name={Aprendizaje autom\'atico (\textit{Machine Learning})}, sort={aprendizaje automatico},
  description={Rama de la inteligencia artificial que desarrolla modelos capaces de aprender patrones a partir de datos, sin ser programados de forma expl\'icita para cada tarea.}}
\newglossaryentry{aprendizajeprofundo}{name={Aprendizaje profundo (\textit{Deep Learning})}, sort={aprendizaje profundo},
  description={Rama del aprendizaje autom\'atico que utiliza redes neuronales de muchas capas para extraer representaciones de los datos de forma autom\'atica.}}
\newglossaryentry{aprendizajesupervisado}{name={Aprendizaje supervisado}, sort={aprendizaje supervisado},
  description={Enfoque de entrenamiento que utiliza datos previamente etiquetados, formados por pares de entrada y salida esperada.}}
\newglossaryentry{aumentodatos}{name={Aumento de datos}, sort={aumento de datos},
  description={T\'ecnica que genera variantes de los datos de entrenamiento para ampliar su diversidad. En este trabajo se aplica la reflexi\'on horizontal de los videos.}}
\newglossaryentry{backend}{name={\textit{Backend}}, sort={backend},
  description={Parte de una aplicaci\'on que se ejecuta en el servidor. Recibe las solicitudes, procesa los datos y se comunica con el modelo de inferencia.}}
\newglossaryentry{balanceo}{name={Balanceo corporal}, sort={balanceo corporal},
  description={Movimiento repetitivo del tronco o de todo el cuerpo, hacia adelante y atr\'as o de un lado a otro. Es una de las tres conductas que identifica el sistema.}}
\newglossaryentry{biomarcador}{name={Biomarcador}, sort={biomarcador},
  description={Caracter\'istica observable y medible que se utiliza como indicador de un proceso biol\'ogico o de una condici\'on de salud.}}
\newglossaryentry{dropout}{name={\textit{Dropout}}, sort={dropout},
  description={T\'ecnica de regularizaci\'on que desactiva al azar una fracci\'on de las neuronas durante el entrenamiento para reducir el sobreajuste.}}
\newglossaryentry{ecolalia}{name={Ecolalia}, sort={ecolalia},
  description={Repetici\'on de palabras o frases dichas por otra persona.}}
\newglossaryentry{endpoint}{name={\textit{Endpoint}}, sort={endpoint},
  description={Direcci\'on de un servicio web a la que se env\'ian solicitudes para ejecutar una operaci\'on concreta.}}
\newglossaryentry{entropia}{name={Entrop\'ia cruzada categ\'orica}, sort={entropia cruzada categorica},
  description={Funci\'on de p\'erdida para clasificaci\'on multiclase que mide la diferencia entre las probabilidades predichas y la clase real.}}
\newglossaryentry{especificidad}{name={Especificidad}, sort={especificidad},
  description={M\'etrica que indica la proporci\'on de casos negativos reales que el modelo clasifica correctamente.}}
\newglossaryentry{estereotipia}{name={Movimiento estereotipado}, sort={movimiento estereotipado},
  description={Movimiento repetitivo y sin una finalidad aparente, como el aleteo de manos o el balanceo corporal.}}
\newglossaryentry{exactitud}{name={Exactitud (\textit{accuracy})}, sort={exactitud},
  description={M\'etrica que indica la proporci\'on de predicciones correctas sobre el total de predicciones.}}
\newglossaryentry{extractor}{name={Extractor de caracter\'isticas}, sort={extractor de caracteristicas},
  description={Parte de una red que transforma una imagen en un vector num\'erico que resume su contenido visual.}}
\newglossaryentry{f1}{name={F1-Score}, sort={f1-score},
  description={M\'etrica que combina la precisi\'on y la sensibilidad mediante su media arm\'onica.}}
\newglossaryentry{fotograma}{name={Fotograma}, sort={fotograma},
  description={Cada una de las im\'agenes fijas que, en secuencia, componen un video.}}
\newglossaryentry{framework}{name={\textit{Framework}}, sort={framework},
  description={Conjunto de herramientas y bibliotecas que proporciona una estructura base para desarrollar software.}}
\newglossaryentry{frontend}{name={\textit{Frontend}}, sort={frontend},
  description={Parte de una aplicaci\'on con la que interact\'ua directamente el usuario, normalmente a trav\'es del navegador.}}
\newglossaryentry{desvanecimiento}{name={Desvanecimiento del gradiente}, sort={desvanecimiento del gradiente},
  description={Problema de las redes profundas o recurrentes en el que el gradiente se reduce al propagarse, lo que dificulta el aprendizaje de dependencias lejanas.}}
\newglossaryentry{hiperparametro}{name={Hiperpar\'ametro}, sort={hiperparametro},
  description={Valor de configuraci\'on que se fija antes del entrenamiento, como la tasa de aprendizaje o el tama\~no de lote.}}
\newglossaryentry{inferencia}{name={Inferencia}, sort={inferencia},
  description={Uso de un modelo ya entrenado para obtener predicciones sobre datos nuevos.}}
\newglossaryentry{matrizconfusion}{name={Matriz de confusi\'on}, sort={matriz de confusion},
  description={Tabla que compara las clases reales con las predichas y muestra los aciertos y errores del modelo para cada clase.}}
\newglossaryentry{neurodesarrollo}{name={Neurodesarrollo}, sort={neurodesarrollo},
  description={Proceso de formaci\'on y maduraci\'on del sistema nervioso, que sustenta el desarrollo cognitivo, motor y social.}}
\newglossaryentry{paradatemprana}{name={Parada temprana (\textit{early stopping})}, sort={parada temprana},
  description={T\'ecnica que interrumpe el entrenamiento cuando la p\'erdida de validaci\'on deja de mejorar, para evitar el sobreajuste.}}
\newglossaryentry{patronmotricidad}{name={Patr\'on de motricidad}, sort={patron de motricidad},
  description={Forma caracter\'istica y repetida de movimiento corporal que puede observarse en una persona.}}
\newglossaryentry{pipeline}{name={\textit{Pipeline}}, sort={pipeline},
  description={Secuencia ordenada de etapas por las que pasan los datos, donde la salida de cada etapa es la entrada de la siguiente.}}
\newglossaryentry{pooling}{name={\textit{Pooling}}, sort={pooling},
  description={Operaci\'on de agrupaci\'on que reduce el tama\~no de los mapas de caracter\'isticas resumiendo los valores de cada regi\'on.}}
\newglossaryentry{precision}{name={Precisi\'on}, sort={precision},
  description={M\'etrica que indica la proporci\'on de predicciones de una clase que son realmente de esa clase.}}
\newglossaryentry{precisionmixta}{name={Precisi\'on mixta}, sort={precision mixta},
  description={T\'ecnica que combina n\'umeros de 16 y 32 bits durante el entrenamiento para reducir el tiempo de c\'omputo y el uso de memoria.}}
\newglossaryentry{red}{name={Red neuronal}, sort={red neuronal},
  description={Modelo computacional formado por capas de unidades conectadas, cuyos par\'ametros se ajustan a partir de los datos durante el entrenamiento.}}
\newglossaryentry{sensibilidad}{name={Sensibilidad (\textit{recall})}, sort={sensibilidad},
  description={M\'etrica que indica la proporci\'on de casos reales de una clase que el modelo identifica correctamente.}}
\newglossaryentry{sobreajuste}{name={Sobreajuste}, sort={sobreajuste},
  description={Situaci\'on en la que un modelo aprende demasiado los detalles de los datos de entrenamiento y pierde capacidad para generalizar a datos nuevos.}}
\newglossaryentry{softmax}{name={\textit{Softmax}}, sort={softmax},
  description={Funci\'on de activaci\'on que convierte las salidas de una red en probabilidades que suman 1, una por cada clase.}}
\newglossaryentry{transferencia}{name={Transferencia de aprendizaje}, sort={transferencia de aprendizaje},
  description={T\'ecnica que reutiliza un modelo entrenado en una tarea para resolver otra, aprovechando el conocimiento ya adquirido.}}
\newglossaryentry{validacioncruzada}{name={Validaci\'on cruzada}, sort={validacion cruzada},
  description={M\'etodo de evaluaci\'on que divide los datos en $k$ pliegues y utiliza cada uno, por turnos, como conjunto de prueba. En este trabajo se emplean cinco pliegues estratificados por clase.}}
\newglossaryentry{ventana}{name={Ventana de conducta}, sort={ventana de conducta},
  description={Intervalo de un video, indicado en las anotaciones, en el que ocurre la conducta de inter\'es.}}
\newglossaryentry{visionartificial}{name={Visi\'on artificial (\textit{Computer Vision})}, sort={vision artificial},
  description={Campo de la inteligencia artificial que busca que los sistemas computacionales adquieran, procesen y comprendan el contenido de im\'agenes y videos.}}
